\clearpage
\fancyhead[L]{\small PALLAS. JUNO.}
\hypertarget{native-274-lower}{}
\begin{center}
PALLAS. JUNO.

\rule{0.6\linewidth}{0.4pt}

Monatliche Correspondenz. Band XIII. S. 313–315. 1806 März.

\rule{0.6\linewidth}{0.4pt}
\end{center}

\par\indent Da uns bis jetzt von der \emph{Pallas} weiter keine Beobachtungen bekannt geworden sind, als die drei auf der Seeberger Sternwarte angestellten, und in dem vorigen Hefte, S. 189 angezeigten, von der \emph{Juno} hingegen noch gar keine, so eilen wir unsern Lesern folgende Nachricht mitzutheilen.

\par\indent Dr.\ \textsc{Gauss} hat die \emph{Pallas} zum erstenmal am 13. Februar wieder gesehen. Da aber diese Beobachtung nur einmal und blos zu dem Zwecke gemacht war, um den Planeten am nächsten heitern Abend durch seine Bewegung aus den übrigen nahe liegenden Sternen, die der Ungewissheit wegen alle mit bemerkt werden mussten, herauszufinden, so hat er uns dieselbe als sehr wenig genau gar nicht geschickt. Am folgenden Abend, den 14. Februar war es ihm nun nicht schwer die fortgerückte Pallas, die sich wie ein Stern 7\textsuperscript{ter} bis 8\textsuperscript{ter} Grösse zeigt, zu erkennen; er verglich sie diesen Abend mit 54 Eridani, so wie an den folgenden Abenden mit Sternen der Hist.\ cél. Hier diese Beobachtungen: [dieselben wie in diesem Bande Seite 267.]

\par\indent Diesen Beobachtungen zu Folge gibt Dr.\ \textsc{Gauss} Ephemeride in der M.\,C.\ XI.\ B.\ S.\ 376 [S.\ 263 d.\ B.] die gerade Aufsteigung dieses Planeten um 27 Minuten, die Abweichung um 5 Minuten zu gross: „Es ist sonderbar“, bemerkt Dr.\ \textsc{Gauss}, „dass erstere Differenz gerade auf die entgegengesetzte Seite fällt, wie nach Ihren Beobachtungen im November vorigen Jahres (siehe vorig. Heft, S. 189); in der Grösse des Fehlers zeigt sich nunmehr, wie ich glaube, der Einfluss der Störungen.“

\par\indent Die \emph{Juno} hat er am 16. Februar zuerst wieder gesehen, und am 17. Februar als \emph{solche} wieder beobachtet. Hier seine Beobachtungen [dieselben wie in diesem Bande Seite 267].

\par\indent Hienach gäbe Dr.\ \textsc{Gauss} Ephemeride (M.\,C.\ XI.\ B.\ S.\ 477 [S.\ 265 d.\ B.]) die gerade Aufsteigung um eine oder anderthalb Minuten zu gross, die Abweichung um eine Minute zu klein, welche, gerade nach einem Jahre Zwischenzeit, eine ganz unbedeutende Abweichung ist, und woraus man zugleich schliessen kann, dass Dr.\ \textsc{Gauss} letzte Beobachtung vom 20. Februar 1805, worauf sich seine V.\ Ele%
\SourceLeaf{275}%
mente gründeten, vorzüglich gut gewesen sein muss. Dr.~\textsc{Gauss} schreibt bei dieser Gelegenheit: „die Juno ist sehr lichtschwach, ich schätze sie etwa einem Stern eilfter Grösse gleich, doch lässt sie sich bei heiterer Luft am Kreis-Micrometer noch gut beobachten, schwerlich wird dies aber am Passagen-Instrumente und Kreis oder Quadranten, wenn man beleuchten muss, möglich sein. Mir werden indess diese Beobachtungen gewöhnlich sehr beschwerlich, da ich nur ein Spiegel-Telescop von \textsc{Short} dazu anzuwenden habe, dessen Gesichtsfeld sehr klein ist, daher es die meiste Zeit an passenden und gut bestimmten Sternen fehlt, die nahe liegen und durch das enge Feld mit durchgehen, der oftmals auch nicht kleinen Schwierigkeit zu gedenken, solche kleine Planeten nur erst aufzufinden, wenn gar keine bekannten Sterne nahe dabei sind, um als Leitsterne zu dienen.“

\par\indent Mit diesen Nachrichten verbinden wir hier noch eine, unsern astronomischen Lesern gewiss höchst angenehme, dass nemlich Dr.~\textsc{Gauss} uns berichtet, dass er sich jetzt hauptsächlich mit verschiedenen einzelnen Materien seiner Methode, die Planetenbahnen zu bestimmen, beschäftige, und in kurzem hoffe, die Ausarbeitung, Vollendung und Herausgabe eines eigenen Werks darüber ernstlich betreiben zu können.

\begin{flushright}
[\textsc{von Zach}.]
\end{flushright}
